\section*{Capítulo \ref{ch:b3:x-ray-diffraction} --- Cristalografia e difração de raios X}

\begin{solution}{exo:b3:x-ray-diffraction:1}
$d = a/\sqrt N$: $d_{111} = \qty{325.66}{pm}$, $d_{200} = \qty{282.03}{pm}$, $d_{220} =
\qty{199.42}{pm}$.
\end{solution}

\begin{solution}{exo:b3:x-ray-diffraction:2}
CFC: as primeiras reflexões são (111) e (200). $\sin\theta = \lambda\sqrt N/2a$: $2\theta =
\ang{43.32}$ e \ang{50.45}.
\end{solution}

\begin{solution}{exo:b3:x-ray-diffraction:3}
P: todas as sete. I ($h + k + l$ par): 110, 200, 211, 220. F (não mistos): 111, 200, 220.
\end{solution}

\begin{solution}{exo:b3:x-ray-diffraction:4}
$\mathbf a^* = \mathbf e_x/a$, $\mathbf b^* = \mathbf e_y/b$: uma rede retangular de lados $1/a$
e $1/b$, em que o lado longo da rede direta se torna o curto.
\end{solution}

\begin{solution}{exo:b3:x-ray-diffraction:5}
Razões de $\sin^2\theta$ $1 : 2 : 3 : 4 : 5$, isto é, $N = 2, 4, 6, 8, 10$: corpo centrado (110,
200, 211, 220, 310). Pela primeira linha, $d_{110} = \lambda/2\sin(\ang{20.135}) =
\qty{223.77}{pm}$ e $a = d\sqrt2 = \qty{316.5}{pm}$: o tungstênio.
\end{solution}

\begin{solution}{exo:b3:x-ray-diffraction:6}
$Z = \rho N_Aa^3/M = 1.99 \times 6.022\times10^{23} \times (6.2929\times10^{-8})^3/74.6 = 4.0$.
\end{solution}

\begin{solution}{exo:b3:x-ray-diffraction:7}
$\beta = \qty{0.40}{\degree} = \qty{6.98e-3}{rad}$, $\cos\theta = \cos\ang{15.85} = 0.962$:
$L = 0.9 \times 0.15406/(6.98\times10^{-3} \times 0.962) = \qty{21}{nm}$.
\end{solution}

\begin{solution}{exo:b3:x-ray-diffraction:8}
$F_{hkl} = f_{\ce{Cs+}} + f_{\ce{Cl-}}(-1)^{h+k+l}$: $F_{100} = 54 - 18 = 36$, $F_{110} = 72$:
intensidades na razão 1 : 4. O centro é ocupado por um íon diferente do dos
vértices: a rede é primitiva (uma rede de corpo centrado exige conteúdos idênticos
nos dois pontos), e 100 está presente.
\end{solution}

\begin{solution}{exo:b3:x-ray-diffraction:9}
Cada átomo tem uma cópia deslocada de $(\frac12,\frac12,0)$, o que multiplica $F$ por $1 +
\eu^{\iu\pi(h+k)} = 1 + (-1)^{h+k}$, nulo para $h + k$ ímpar.
\end{solution}

\begin{solution}{exo:b3:x-ray-diffraction:10}
Os quasicristais têm ordem de longo alcance sem periodicidade: não têm rede,
de modo que a restrição, que diz respeito às redes, não se aplica a eles. Suas
manchas de difração nítidas levaram a definir um cristal como todo sólido de
difratograma essencialmente discreto, periódico ou não.
\end{solution}

\begin{solution}{exo:b3:x-ray-diffraction:11}
O deslizamento $c$ leva $(x, y, z)$ a $(x, \frac12 - y, \frac12 + z)$. Para $k = 0$, os dois
fatores de fase são $\eu^{2\pi\iu(hx + lz)}$ e $\eu^{2\pi\iu(hx + lz + l/2)} = (-1)^l
\eu^{2\pi\iu(hx+lz)}$: sua soma se anula para $l$ ímpar.
\end{solution}

\begin{solution}{exo:b3:x-ray-diffraction:12}
Uma inversão, um espelho ou um deslizamento transformam uma molécula em sua
imagem especular, de modo que um cristal que contivesse um deles conteria os dois
enantiômeros. Um enantiômero puro só cristaliza nos 65 \omterm{def:b3:x-ray-diffraction:space-group}{grupos espaciais} construídos
com rotações, \omterm{def:b3:x-ray-diffraction:space-group}{eixos helicoidais} e translações (grupos de Sohncke), como
\textit{P}$2_12_12_1$. Pela lei de Friedel, o difratograma de um enantiômero é igual ao
do outro; a configuração absoluta é obtida dos pequenos desvios causados pelo
espalhamento anômalo dos átomos mais pesados.
\end{solution}

\begin{solution}{pb:b3:x-ray-diffraction:1}
\textbf{1.} $\theta = \ang{14.171}$, $d = \lambda/2\sin\theta = \qty{314.64}{pm}$.
\textbf{2.} 314.64, 222.49, 181.66, 157.32, 140.71, \qty{128.45}{pm}.
\textbf{3.} 1.000, 2.000, 3.000, 4.000, 5.000, 6.000.
\textbf{4.} 1, 2, 3, 4, 5, 6.
\textbf{5.} Sim, como (100), (110), (111), (200), (210), (211) de uma célula cúbica primitiva
com $a' = \qty{314.6}{pm}$.
\textbf{6.} $N = 7$ (e 15, 23\dots), que não é soma de três quadrados; a primeira
diferença apareceria na oitava linha. Seis linhas não permitem distinguir uma
célula P de aresta $a'$ de uma célula F de aresta $2a'$.
\textbf{7.} $N = 4, 8, 12, 16, 20, 24$: (200), (220), (222), (400), (420), (422).
\textbf{8.} $a = 2d_{200} = \qty{629.3}{pm}$.
\textbf{9.} \ce{KCl} (\qty{629.29}{pm}).
\textbf{10.} $F = 4[f(\ce{K+}) - f(\ce{Cl-})]$ para índices todos ímpares, e os dois íons têm 18
elétrons cada um: as amplitudes se cancelam.
\textbf{11.} Sim, para os dois: \ce{Na+} (10) e \ce{Cl-} (18), \ce{K+} (18) e \ce{Br-} (36)
espalham de modo diferente; a (111) do \ce{NaCl} estaria em \ang{27.36} e a do \ce{KBr} em
\ang{23.38}, ambas dentro do intervalo registrado.
\textbf{12.} Os raios X veem \ce{K+} e \ce{Cl-} como quase idênticos; íons idênticos em
um arranjo de sal-gema formam um arranjo cúbico simples de aresta $a/2$, uma célula aparente.
\textbf{13.} Quatro \ce{KCl}.
\textbf{14.} $\rho = 4 \times 74.6/(6.022\times10^{23} \times (6.293\times10^{-8})^3) =
\qty{1.99}{g/cm^3}$.
\textbf{15.} Seis e seis (octaedros do outro íon).
\textbf{16.} $a/2 = \qty{314.6}{pm}$.
\textbf{17.} (440), $N = 32$, em \ang{87.65} (estando ausentes as todas ímpares (333)/(511)).
\textbf{18.} As intensidades mudam com a orientação preferencial dos cristalitos, a
absorção, o decréscimo dos fatores de espalhamento e a agitação térmica; as
posições dependem apenas da rede.
\textbf{19.} De $d = \lambda/2\sin\theta$, $\Delta d/d = -\cot\theta\,\Delta\theta$: o mesmo erro
angular dá um erro relativo menor para $\theta$ grande.
\textbf{20.} $d_{420} = \lambda/2\sin\ang{33.190} = \qty{140.71}{pm}$, $a = d\sqrt{20} =
\qty{629.3}{pm}$.
\textbf{21.} $L = 0.9 \times 0.15406/(4.36\times10^{-3} \times \cos\ang{33.19}) = \qty{38}{nm}$.
\textbf{22.} Não: $\beta = \ang{0.0095}$, muito abaixo de uma largura instrumental típica.
\textbf{23.} \textbf{$a = \qty{629.3}{pm}$: o pó é cloreto de potássio}, com as reflexões
ímpares silenciadas por dois íons com o mesmo número de elétrons.
\end{solution}
